\section{Theoretical insights}
\label{sec:theory}
% Under the setting described in \Cref{subsec:notations}, we detail our new theoretical analysis of WA, introduced in \Cref{sec:limithessian}. % theoretically analyze 
% It relies on a new bias-variance-covariance-locality decomposition of its expected \ood generalization error, derived in \Cref{subsec:biasvariance}.
% This decomposition highlights the importance of diversity across averaged models.
% %In \Cref{subsec:biasvariance}, we derive the expected \ood generalization error of WA via a new bias-variance-covariance-locality decomposition.
% We clarify in \Cref{subsec:analysisbvc} the conditions of WA's \ood success by showing the close links between bias / correlation shift and variance / diversity shift.
% We then present our two main results: the bias term dominates in the decomposition under correlation shift (see \Cref{subsec:expression_ood_bias}) and variance dominates under diversity shift (see \Cref{subsec:expression_ood_var}).
% %Based on this analysis, we propose to increase the diversity across models averaged in weights to tackle distribution shift.
% % This analysis provides guidelines to understand when WA succeeds or fails.
% This analysis provides guidelines to understand when WA succeeds, and how to improve it.
% Moreover, this decomposition highlights the importance of diversity across averaged models.
% %In \Cref{subsec:biasvariance}, we derive the expected \ood generalization error of WA via a new bias-variance-covariance-locality decomposition.
%  by showing the close links between bias / correlation shift and variance / diversity shift.
%Based on this analysis, we propose to increase the diversity across models averaged in weights to tackle distribution shift.
% This will enable us to undertand when WA 
% In particular, we will show that the bias term dominates in the decomposition under correlation shift (see \Cref{subsec:expression_ood_bias}) and variance dominates under diversity shift (see \Cref{subsec:expression_ood_var}).
Under the setting described in \Cref{subsec:notations}, we introduce WA in \Cref{sec:limithessian} and decompose its expected \ood error in \Cref{subsec:biasvariance}.
Then, we separately consider the four terms of this bias-variance-covariance-locality decomposition in \Cref{subsec:analysisbvc}.
This theoretical analysis will allow us to better understand when WA succeeds, and most importantly, how to improve it empirically in \Cref{sect:dwa}.%
\subsection{Notations and problem definition}%
\label{subsec:notations}
% \vspace{-0.6em}%
\paragraph{Notations.}%
We denote $\X$ the input space of images, $\Y$ the label space
and $\ell:\Y^2 \rightarrow \mathbb{R}_+$ a loss function.
$S$ is the training (source) domain with distribution $p_S$, and $T$ is the test (target) domain with distribution $p_T$.
For simplicity, we will indistinctly use the notations $p_S$ and $p_T$ to refer to the joint, posterior and marginal distributions of $(X, Y)$.
% $p_S$ and $p_T$ will indistinctly refer to the joint, posterior and marginal distributions, respectively $p(X,Y), p(Y|X), p(X)$.
We note $f_S, f_T:\X\rightarrow\Y$ the source and target labeling functions.
We assume that there is no noise in the data: then $f_S$ is defined on $\X_S \triangleq \{x\in \X / p_S(x)>0\}$ by $\forall (x, y) \sim p_S,f_S(x)=y$ and similarly $f_T$ is defined on $\X_T \triangleq \{x\in \X / p_T(x)>0\}$ by $\forall (x, y) \sim p_T, f_T(x)=y$.%
% \vspace{-0.1em}%
\paragraph{Problem.}
We consider a neural network (NN) $f(\cdot, \theta):\X\rightarrow\Y$ made of a fixed architecture $f$ with weights $\theta$.
We seek $\theta$ minimizing the target generalization error:
\begin{equation} \label{eq:gen_error}
    \mathcal{E}_T(\theta)= \mathbb{E}_{(x,y) \sim p_T}[\ell(f(x, \theta),y)].
\end{equation}
% $\mathcal{E}_T(\theta)= \mathbb{E}_{(x,y) \sim p_T}[\ell(f(x, \theta),y)]$.
$f(\cdot, \theta)$ should approximate $f_T$ on $\X_T$. However, this is complex in the \ood setup because we only have data from domain $S$ in training, related yet different from $T$.
The differences between $S$ and $T$ are due to distribution shifts (\ie the fact that
$p_S(X,Y) \neq p_T(X,Y)$) which are decomposed per \cite{ye2021odbench} into
\textbf{diversity shift} (\aka covariate shift), when marginal distributions differ (\ie $p_S(X) \neq p_T(X)$), and
\textbf{correlation shift} (\aka concept shift), when posterior distributions differ (\ie $p_S(Y|X) \neq p_T(Y|X)$ and $f_S \neq f_T$).
The weights are typically learned on a training dataset $d_S$ from $S$ (composed of $n_S$ i.i.d. samples from $p_S(X,Y)$) with a configuration $c$, which contains all other sources of randomness in learning (\eg initialization, hyperparameters, training stochasticity, epochs, \etc). %number of steps, \etc).
We call $l_S=\{d_S, c\}$ a learning procedure on domain $S$, and explicitly write $\theta(l_S)$ to refer to the weights obtained after stochastic minimization of $1/n_S \sum_{(x,y)\in d_S} \ell(f(x, \theta), y)$ \wrt $\theta$ under $l_S$.
% \vspace{-0.1em}%
\subsection{Weight averaging for OOD and limitations of current analysis}%
\label{sec:limithessian}%
\paragraph{Weight averaging.}%
We study the benefits of combining $M$ individual member weights $\{\theta_m\}_{m=1}^M \triangleq \{\theta(l_S^{(m)})\}_{m=1}^M$ obtained from $M$ (potentially correlated) identically distributed (i.d.) learning procedures $L_S^M\triangleq\{l_S^{(m)}\}_{m=1}^M$.
Under conditions discussed in \Cref{subsec:sharedinithpws}, these $M$ weights can be averaged despite nonlinearities in the architecture $f$.
Weight averaging (WA) \cite{izmailov2018}, defined as:
\begin{equation} \label{eq:f_wa}
    f_{\text{WA}} \triangleq f(\cdot, \theta_{\text{WA}}), \text{~where~} \theta_{\text{WA}}\triangleq\theta_{\text{WA}}(L_S^M) \triangleq 1/M \sum\nolimits_{m=1}^{M} \theta_m,%
\end{equation}
is the state of the art \cite{cha2021wad,arpit2021ensemble} on DomainBed \cite{gulrajani2021in} when the weights $\{\theta_m\}_{m=1}^M$ are sampled along a single training trajectory (a description we refine in \Cref{remark:identical_assumption} from \Cref{app:proof_bvc}).
\paragraph{Limitations of the flatness-based analysis.}
To explain this success, Cha \textit{et al.} \cite{cha2021wad} argue that flat minima generalize better; indeed, WA flattens the loss landscape.
Yet, as shown~in~\Cref{app:limit_proof_swad}, this analysis does not fully explain WA's spectacular results on DomainBed.
First, flatness does not act on distribution shifts thus the \ood error is uncontrolled with their upper bound (see \Cref{app:flatness_distribution_shifts}).
Second, this analysis does not clarify why WA outperforms \reb{Sharpness-Aware Minimizer (SAM) \cite{foret2021} for \ood generalization, even though SAM} directly optimizes flatness (see \Cref{app:mav_better_than_sam}).
Finally, it does not justify why combining WA and SAM succeeds in \iid \cite{questionsflatminima22} yet fails in \ood (see \Cref{app:mav_sam_failure}).
These observations motivate a new analysis of WA; we propose one below that better explains these~results. %
\subsection{Bias-variance-covariance-locality decomposition}
\label{subsec:biasvariance}
We now introduce our bias-variance-covariance-locality decomposition which extends the bias-variance decomposition \cite{kohavi1996bias} to WA.
In the rest of this theoretical section, $\ell$ is the Mean Squared Error for simplicity: yet, our results may be extended to other losses as in \cite{domingos2000unified}.
In this case, the expected error of a model with weights $\theta(l_S)$ \wrt the learning procedure $l_S$ was decomposed in \cite{kohavi1996bias} into:
\begin{equation} \tag{BV}\label{eq:b_v}
    \mathbb{E}_{l_S}\mathcal{E}_T(\theta(l_S)) = \mathbb{E}_{(x,y)\sim p_T}[\biasb^2(x, y)+\varb(x)],
\end{equation}
where $\biasb(x,y), \varb(x)$ are the bias and variance of the considered model \wrt a sample $(x,y)$, defined later in \Cref{eq:b_var_cov}.
To decompose WA's error, we leverage the similarity (already highlighted in \cite{izmailov2018}) between WA and functional ensembling (ENS) \cite{Lakshminarayanan2017,dietterich2000ensemble}, a more traditional way to combine a collection of weights.
More precisely, ENS averages the predictions, $f_{\text{ENS}} \triangleq f_{\text{ENS}}(\cdot, \{\theta_m\}_{m=1}^M) \triangleq 1/M \sum_{m=1}^{M} f(\cdot, \theta_m)$.
\Cref{lemma:wa_ensembling} establishes that $f_{\text{WA}}$ is a first-order approximation of $f_{\text{ENS}}$ when $\{\theta_m\}_{m=1}^M$ are close in the weight space.
\begin{lemma}[WA and ENS. Proof in \Cref{app:wa_loss}. Adapted from \cite{izmailov2018,Wortsman2022ModelSA}.] \label{lemma:wa_ensembling}
    Given $\{\theta_m\}_{m=1}^M$ with learning procedures $L_S^M\triangleq\{l_S^{(m)}\}_{m=1}^M$. Denoting $\Delta_{L_S^M}=\max_{m=1}^{M}\left\|\theta_m-\theta_{\text{WA}}\right\|_2$, $\forall (x,y) \in \X \times \Y$:%
    \begin{equation*}
        f_{\text{WA}}(x) = f_{\text{ENS}}(x) + O(\Delta^2_{L_S^M}) \text{~and~} \ell\left(f_{\text{WA}}(x), y\right) = \ell\left(f_{\text{ENS}}(x) , y\right) + O(\Delta_{L_S^M}^2).%
    \end{equation*}%
\end{lemma}%
This similarity is useful since \Cref{eq:b_v} was extended into a bias-variance-covariance decomposition for ENS in \cite{ueda1996generalization,brown2005between}.
We can then derive the following decomposition of WA's expected test error. To take into account the $M$ averaged weights, the expectation is over the joint distribution describing the $M$ identically distributed (i.d.) learning procedures $L_S^M\triangleq\{l_S^{(m)}\}_{m=1}^M$.
\begin{proposition}[Bias-variance-covariance-locality decomposition of the expected generalization error of WA in \ood. Proof in \Cref{app:proof_bvc}.]
    \label{prop:b_var_cov}
    Denoting $\bar{f}_S\left(x\right) = \mathbb{E}_{l_S} \left[f\left(x,\theta\left(l_S\right)\right)\right]$, under identically distributed learning procedures $L_S^M\triangleq\{l_S^{(m)}\}_{m=1}^M$, the expected generalization error on domain $T$ of $\theta_{\text{WA}}(L_S^M)\triangleq\frac{1}{M} \sum\nolimits_{m=1}^{M} \theta_m$ over the joint distribution of $L_S^M$ is:%
    \begin{equation}
        \begin{aligned}
             \mathbb{E}_{L_S^M}\mathcal{E}_T(\theta_{\text{WA}}(L_S^M)) &= \mathbb{E}_{(x,y)\sim p_T}\Big[\biasb^2(x, y)+\frac{1}{M} \varb(x)+\frac{M-1}{M} \covb(x)\Big] + O(\Deltab^2), \\
             \text{where~}\biasb(x,y)&=y-\bar{f}_S\left(x\right), \\
             \text{and~}\varb(x)&=\mathbb{E}_{l_S}\left[\left(f(x, \theta(l_S)) - \bar{f}_S\left(x\right)\right)^{2}\right], \\
             \text{and~}\covb(x)&= \mathbb{E}_{l_S,l_S'}\left[\left(f(x,\theta(l_S))-\bar{f}_S\left(x\right)\right)\left(f(x,\theta(l_S')))-\bar{f}_S\left(x\right)\right)\right], \\
             \text{and~}\Deltab^2&=\mathbb{E}_{L_S^M}\Delta_{L_S^M}^2 \text{~with~} \Delta_{L_S^M}=\max_{m=1}^{M}\left\|\theta_m-\theta_{\text{WA}}\right\|_2.%
        \end{aligned} \tag{BVCL}\label{eq:b_var_cov}%
    \end{equation}%
    $\covb$ is the prediction covariance between two member models whose weights are averaged.
    The locality term $\Deltab^2$ is the expected squared maximum distance between weights and their average.
\end{proposition}
\Cref{eq:b_var_cov} decomposes the \ood error of WA into four terms.
The bias is the same as that of each of its i.d. members.
WA's variance is split into the variance of each of its i.d. members divided by $M$ and a covariance term.
The last locality term constrains the weights to ensure the validity of our approximation.
In conclusion, combining $M$ models divides the variance by $M$ but introduces the covariance and locality terms which should be controlled along bias to guarantee low \ood error.
\subsection{Analysis of the bias-variance-covariance-locality decomposition}%
\label{subsec:analysisbvc}
We now analyze the four terms in \Cref{eq:b_var_cov}.
We show that bias dominates under correlation shift (\Cref{subsec:expression_ood_bias}) and variance dominates under diversity shift (\Cref{subsec:expression_ood_var}).
Then, we discuss a trade-off between covariance, reduced with diverse models (\Cref{subsec:expression_cov_div}), and the locality term, reduced when weights are similar (\Cref{subsec:expression_loc_lir}).
This analysis shows that \textit{WA is effective against diversity shift when $M$ is large and when its members are diverse but close in the weight space}.
\subsubsection{Bias and correlation shift (and support mismatch)}
\label{subsec:expression_ood_bias}
We relate \ood bias to correlation shift \cite{ye2021odbench} %, \ie posterior shift between $p_S(Y|X)$ and $p_T(Y|X)$.
under \Cref{ass:no_bias_iid}, where $\bar{f}_{S}\left(x\right) \triangleq \mathbb{E}_{l_S} \left[f\left(x,\theta\left(l_S\right)\right)\right]$.
As discussed in \Cref{app:bias_correlation_ass_noiidbias}, \Cref{ass:no_bias_iid} is reasonable for a large NN trained on a large dataset representative of the source domain $S$. 
It is relaxed in \Cref{prop:app_bias_full} from \Cref{app:bias_correlation}.
\begin{assumption}[Small \iid bias]%
    $\exists \epsilon > 0 \text{~small~s.t.~} \forall x\in \X_S, |f_{S}\left(x\right)-\bar{f}_{S}\left(x\right)|\leq \epsilon$.%
    \label{ass:no_bias_iid}%
    % \vspace{-0.5em}%
\end{assumption}%
\begin{proposition}[\ood bias and correlation shift. Proof in \Cref{app:bias_correlation}] \label{prop:bias}
    With a bounded difference between the labeling functions $f_T-f_S$ on $\X_T \cap \X_S$, under \Cref{ass:no_bias_iid}, the bias on domain $T$ is:
    \begin{equation}\label{eq:bias}
        \begin{aligned}
            \mathbb{E}_{(x,y)\sim p_T}[\biasb^2(x, y)] &= \text{Correlation shift} + \text{Support mismatch} + O(\epsilon), \\
            \text{where~} \text{Correlation shift} &= \int_{\X_T \cap \X_S} \left(f_T(x)-f_S(x)\right)^{2} p_T(x) dx, \\
            \text{and~} \text{Support mismatch} &= \int_{\X_T \setminus \X_S} \left(f_T(x)-\bar{f}_{S}\left(x\right)\right)^{2} p_T(x) dx.%
        \end{aligned}%
    \end{equation}
\end{proposition}%
We analyze the first term by noting that $f_T(x) \triangleq \mathbb{E}_{p_T}[Y|X=x]$ and $f_S(x) \triangleq \mathbb{E}_{p_S}[Y|X=x],$ $\forall x \in \X_T \cap \X_S$.
This expression confirms that our correlation shift term measures shifts in posterior distributions between source and target, as in \cite{ye2021odbench}.
It increases in presence of spurious correlations: \eg on ColoredMNIST \cite{arjovsky2019invariant} where the color/label correlation is reversed at test time.
The second term is caused by support mismatch between source and target.
It was analyzed in \cite{ruan2022optimal} and shown irreducible in their \enquote{No free lunch for learning representations for DG}.
Yet, this term can be tackled if we transpose the analysis in the feature space rather than the input space. This motivates encoding the source and target domains into a shared latent space, \eg by pretraining the encoder on a task with minimal domain-specific information as in \cite{ruan2022optimal}.

This analysis explains why WA fails under correlation shift, as shown on ColoredMNIST in \Cref{app:failure_corr_shift}.
Indeed, combining different models does \textit{not} reduce the bias.
\Cref{subsec:expression_ood_var} explains that WA is however efficient against diversity shift.%
\subsubsection{Variance and diversity shift}
\label{subsec:expression_ood_var}
Variance is known to be large in \ood \cite{d2020underspecification} and to cause a phenomenon named  underspecification, when models behave differently in \ood despite similar test \iid accuracy.
We now relate \ood variance to diversity shift \cite{ye2021odbench} in a simplified setting.
We fix the source dataset $d_S$ (with input support $X_{d_S}$), the target dataset $d_T$ (with input support $X_{d_T}$) and the network's initialization.
We get a closed-form expression for the variance of $f$ over all other sources of randomness under \Cref{ass:infinite_width,ass:inter_sample}.
\begin{assumption}[Kernel regime] $f$ is in the kernel regime \cite{Jacot2018,daniely2017sgd}.%
\label{ass:infinite_width}%
% \vspace{-0.5em}%
\end{assumption}%
This states that $f$ behaves as a Gaussian process (GP); it is reasonable if $f$ is a wide network \cite{Jacot2018,Lee2018DeepNN}.
The corresponding kernel $K$ is the neural tangent kernel (NTK) \cite{Jacot2018} depending only on the initialization.
GPs are useful because their variances have a closed-form expression (\Cref{app:nns_as_gps}).
To simplify the expression of variance, we now make \Cref{ass:inter_sample}.
\begin{assumption}[Constant norm and low intra-sample similarity on $d_S$]
    % $\exists 0\leq \epsilon \ll \lambda_S \text{~s.t.~} \forall (x_S, x_S^\prime \neq x_S) \in X_{d_S}^2,|K(x_S,x_S^\prime)|\leq \epsilon$.%
    $\exists (\lambda_S, \epsilon) \text{~with~}\linebreak 0 \leq \epsilon \ll \lambda_S \text{~such~that~} \forall x_S \in X_{d_S}, K(x_S,x_S)=\lambda_S \text{~and~} \forall x_S^\prime \neq x_S \in X_{d_S},|K(x_S,x_S^\prime)|\leq \epsilon$.%
    \label{ass:inter_sample}%
    % \vspace{-0.5em}%
\end{assumption}%
This states that training samples have the same norm (following standard practice \cite{Lee2018DeepNN,ah2010normalized,ghojogh2021reproducing,rennie2005}) and weakly interact \cite{He2020The,seleznova2022neural}. This assumption is further discussed and relaxed in \Cref{app:discussion_inter_sample}.
We are now in a position to relate variance and diversity shift when~$\epsilon \to 0$.%
\clearpage
\begin{proposition}[\ood variance and diversity shift. Proof in \Cref{app:var_diversity}] \label{prop:var}%
    Given $f$ trained on source dataset $d_S$ (of size $n_S$) with NTK $K$, under \Cref{ass:infinite_width,ass:inter_sample}, the variance on dataset $d_T$ is:%
    \begin{equation} \label{eq:int_var}%
        \mathbb{E}_{x_T\in X_{d_T}}[\varb(x_T)] = \frac{n_S}{2\lambda_S}\text{MMD}^{2}(X_{d_S}, X_{d_T}) + \lambda_T - \frac{n_S}{2\lambda_S}\beta_T + O(\epsilon),%
    \end{equation}%
    where $\text{MMD}$ is the empirical Maximum Mean Discrepancy in the RKHS of $K^2(x,y)=(K(x,y))^2$;$\lambda_T \triangleq \mathbb{E}_{x_T \in X_{d_T}} K\left(x_T, x_T\right)$ and $\beta_T \triangleq \mathbb{E}_{(x_T,x_T^\prime)\in X_{d_T}^2, x_T \neq x_T^\prime} K^2\left(x_T, x_T^{\prime}\right)$ are the empirical mean similarities respectively measured between identical (\wrt $K$) and different (\wrt $K^2$) samples averaged over $X_{d_T}$.%
\end{proposition}
The MMD empirically estimates shifts in input marginals, \ie between $p_S(X)$ and $p_T(X)$. Our expression of variance is thus similar to the diversity shift formula in \cite{ye2021odbench}: MMD replaces the $L_1$ divergence used in \cite{ye2021odbench}.
The other terms, $\lambda_T$ and $\beta_T$, both involve internal dependencies on the target dataset $d_T$: they are constants \wrt $X_{d_T}$ and do not depend on distribution shifts.
At fixed $d_T$ and under our assumptions, \Cref{eq:int_var} shows that variance on $d_T$ decreases when $X_{d_S}$ and $X_{d_T}$ are closer (for the MMD distance defined by the kernel $K^2$) and increases when they deviate. Intuitively, the further $X_{d_T}$ is from $X_{d_S}$, the less the model's predictions on $X_{d_T}$ are constrained after fitting $d_S$.

This analysis shows that WA reduces the impact of diversity shift as combining $M$ models divides the variance per $M$. This is a strong property achieved \textit{without requiring data from the target domain}.
\subsubsection{Covariance and diversity}
\label{subsec:expression_cov_div}
The covariance term increases when the predictions of $\{f(\cdot, \theta_m)\}_{m=1}^M$ are correlated.
In the worst case where all predictions are identical, covariance equals variance and WA is no longer beneficial.
On the other hand, the lower the covariance, the greater the gain of WA over its members; this is derived by comparing
\Cref{eq:b_v,eq:b_var_cov}, as detailed in \Cref{app:proof_wa_ind}.
It motivates tackling covariance by encouraging members to make different predictions, thus to be functionally diverse.
Diversity is a widely analyzed concept in the ensemble literature \cite{Lakshminarayanan2017}, for which numerous measures have been introduced \cite{kuncheva2003measures,aksela2003comparison,DBLP:journals/corr/abs-1905-00414}.
\reb{In \Cref{sect:dwa}, we aim at decorrelating the learning procedures to increase members' diversity and reduce the covariance term}.

% Deep ensembles \cite{Lakshminarayanan2017} usually seek zero covariance by members' independence, \reb{\eg via explicit adversarial regularization in \cite{rame2021dice}: we implicitly follow the same goal when s in \Cref{sect:dwa}}.
% we follow this objective in \Cref{sect:dwa} by decorrelating the learning procedures.
% For example, \cite{rame2021dice} explicitly sought zero covariance by member's independence in Deep ensembles \cite{Lakshminarayanan2017}.
\subsubsection{Locality and linear mode connectivity}
\label{subsec:expression_loc_lir}
To ensure that WA approximates ENS, the last locality term $O(\Deltab^2)$ constrains the weights to be close.
Yet, the covariance term analyzed in \Cref{subsec:expression_cov_div} is antagonistic, as it motivates functionally diverse models.
Overall, to reduce WA's error in \ood, we thus seek a good trade-off between diversity and locality.
In practice, we consider that the main goal of this locality term is to ensure that the weights are averageable despite the nonlinearities in the NN such that WA's error does not explode.
This is why in \Cref{sect:dwa}, we empirically relax this locality constraint and simply require that the weights are linearly connectable in the loss landscape, as in the linear mode connectivity \cite{Frankle2020}. We empirically verify later in \Cref{fig:home0_samediffruns_net_soup} that the approximation $f_{\text{WA}} \approx f_{\text{ENS}}$ remains valid even in this case.
